% ECONOMICS_PROBLEM: {"id": "C63-209", "title": "Exact complexity of an Arrow--Debreu equilibrium", "jel_code": "C63", "source_id": "OP-0573"}
% !TeX program = xelatex
\documentclass[11pt,a4paper]{article}
\usepackage[margin=23mm]{geometry}
\usepackage{fontspec}
\setmainfont{Latin Modern Roman}
\usepackage{amsmath,amssymb,mathtools}
\usepackage{xcolor,enumitem,booktabs,longtable,array,xurl,hyperref}
\definecolor{muted}{HTML}{53636C}
\hypersetup{colorlinks=true,urlcolor=blue,linkcolor=blue}
\setlength{\parindent}{0pt}
\setlength{\parskip}{5pt}
\newcommand{\R}{\mathbb R}
\newcommand{\E}{\mathbb E}
\newcommand{\N}{\mathbb N}
\newcommand{\ind}{\mathbf 1}
\newcommand{\Var}{\operatorname{Var}}
\newcommand{\Cov}{\operatorname{Cov}}
\newcommand{\supp}{\operatorname{supp}}
\DeclareMathOperator*{\argmin}{arg\,min}
\DeclareMathOperator*{\argmax}{arg\,max}
\newcommand{\entrymeta}[3]{{\sffamily\small\color{muted}\textbf{\texttt{#1}}\quad|\quad #2\par #3\par}}
\newcommand{\Problemref}[1]{\texttt{#1}}
\newcommand{\JELref}[2]{\texttt{#2} (JEL \texttt{#1})}
\begin{document}
\section*{Exact complexity of an Arrow--Debreu equilibrium}
\textbf{Problem ID:} \texttt{C63-209}\quad \textbf{JEL:} \texttt{C63}\par
% BEGIN ECONOMICS STATEMENT
\label{problem:OP-0573}
\label{beyond:GE-01}
\entrymeta{GE-01}{General equilibrium and computation}{Published question, open in the cited version}
\subsection*{Definitions and assumptions}
An economy has $L$ commodities, consumers $i=1,\ldots,m$, and firms $j=1,\ldots,n$. Consumer $i$ has a closed convex consumption set $X_i\subset\R^L$ bounded below, a continuous concave nonsatiated utility $u_i:X_i\to\R$, endowment $\omega_i$, and nonnegative profit shares $\alpha_{ij}$ with $\sum_i\alpha_{ij}=1$. Require a bundle $\underline x_i\in X_i$ with $\underline x_i<\omega_i$ coordinatewise. Each closed convex production set $Y_j$ contains $0$; with $Y=\sum_jY_j$, require $Y\cap\R^L_+=Y\cap(-Y)=\{0\}$.

Use the finite input representation of Section~6.1.1 of the reference: rational endowments, shares and consumption lower bounds; convex inequalities and linear equalities satisfying Slater's condition for feasible sets; algebraic pseudogates for the required utility supergradients and constraint subgradients; and an explicit rational bound on feasible production tuples. The promises in that section and its existence theorem are part of the input specification.

An equilibrium consists of $p\in\R^L_+$, $\sum_\ell p_\ell=1$, consumption $x_i$, and production $y_j$ satisfying
\[
 y_j\in\argmax_{y\in Y_j}p\cdot y,\qquad
 x_i\in\argmax_{x\in X_i}\{u_i(x):p\cdot x\le p\cdot\omega_i+\textstyle\sum_j\alpha_{ij}p\cdot y_j\},
\]
\[
 z:=\sum_ix_i-\sum_i\omega_i-\sum_jy_j\le0,\qquad p\cdot z=0.
\]
Here $z\le0$ allows disposal of goods priced at zero.
\subsection*{Formal question}
Is the total search problem of finding such an exact equilibrium FIXP-hard under the polynomial-time reductions used in the reference? FIXP is the class of real search problems reducible to finding an exact fixed point of a continuous algebraic-circuit self-map of a compact convex domain; solutions may be irrational. Together with the cited membership theorem, hardness would give FIXP-completeness for this promised input class.
\subsection*{Scope and status}
The encoding and existence promises are essential: hardness for a market family allowing nonexistence does not automatically establish hardness for this total search problem. The question concerns price-taking competitive allocations and is distinct from weak approximate equilibrium computation. It is explicitly open in the April~2023 source; no universal current-status claim is made.
\subsection*{References for the source of the problem}
Aris Filos-Ratsikas, Kristoffer Arnsfelt Hansen, Kasper H\o gh, and Alexandros Hollender, \emph{FIXP-membership via Convex Optimization: Games, Cakes, and Markets}, arXiv:2111.06878v3, revised 25 April 2023. Definitions and encoding: Sections~6.1 and~6.1.1; open question: Section~7, pp.~44--45. \url{https://arxiv.org/pdf/2111.06878v3}. Journal DOI: \url{https://doi.org/10.1137/22M1472656}.

\subsection*{Common conventions: Source mathematical conventions}

\textbf{Well-defined objects.}
All probability spaces and measurable variables are specified or inherited
from a local statistical model. Expectations used as finite targets require
the stated integrability. A decision rule or estimator is measurable with
respect to its available information; randomized rules may use an auxiliary
independent random variable. Norms without a subscript are Euclidean in
finite-dimensional spaces. An optimizer is requested only when attainment
is assured; otherwise the question concerns the optimal value and
$\varepsilon$-optimal admissible rules for every $\varepsilon>0$.

\textbf{Identification.}
A structural model $m\in\mathcal M$ implies an observable law
$P_m^{\rm obs}$ and target $\theta(m)$. The target is point identified on
$\mathcal M$ if
\[
 P_{m_1}^{\rm obs}=P_{m_2}^{\rm obs}
 \quad\Longrightarrow\quad\theta(m_1)=\theta(m_2)
 \qquad(m_1,m_2\in\mathcal M).
\]
When this fails, the sharp identified set is
\[
 \Theta(P;\mathcal M)=\{\theta(m):m\in\mathcal M,
                                  P_m^{\rm obs}=P\}.
\]
Specifying an intervention or estimand does not establish identification.
Positivity, exclusion, causal sufficiency, measurement restrictions, and
transport assumptions are substantive local inputs, never automatic
consequences of notation.

\textbf{Minimax risk and nontrivial inference.}
For a statistical experiment with data law $P^{(n)}$, target $\theta(P)$,
and nonnegative loss $\ell$, define
\[
 \mathcal R_n(\mathcal P,\ell)
   =\inf_{\widehat\theta_n}\sup_{P\in\mathcal P}
      \E_{P^{(n)}}\ell(\widehat\theta_n,\theta(P)).
\]
The infimum is over measurable estimators. A sharp rate is a sequence
$r_n$ for which $0<c\leq C<\infty$, independent of $n$, satisfies
$c r_n\leq\mathcal R_n\leq C r_n$ for all sufficiently large $n$.
Squared-error parametric risk is of order $n^{-1}$; the corresponding
root-mean-squared error is of order $n^{-1/2}$. These different conventions
are distinguished locally. A uniform asymptotic confidence region $C_n$
must satisfy
\[
 \liminf_{n\to\infty}\inf_{P\in\mathcal P}
     \Pr_{P^{(n)}}\{\theta(P)\in C_n\}\geq 1-\alpha,
 \qquad 0<\alpha<1,
\]
together with an informative diameter or expected-width requirement.
Coverage alone permits the vacuous whole parameter space. Testing questions
similarly impose both size control and power against a declared separated
alternative class.

\textbf{Binary causal baseline.}
When an entry explicitly invokes this baseline, one observes independent
copies of $O=(X,W,Y)$, with $W\in\{0,1\}$ and potential outcomes $Y(0),Y(1)$.
Assume consistency $Y=Y(W)$, no interference, conditional exchangeability
$(Y(0),Y(1))\perp W\mid X$, and overlap
$\epsilon\leq e(X)\leq1-\epsilon$ almost surely for a fixed
$0<\epsilon<1/2$, where
\[
 e(x)=\Pr(W=1\mid X=x),\qquad
 m_w(x)=\E[Y\mid W=w,X=x],\qquad
 \tau(x)=m_1(x)-m_0(x).
\]
These assumptions identify conditional mean effects almost everywhere
under the covariate law. Pointwise or uniform effects require appropriate
regularity and a specified support. Continuous treatment, longitudinal
data, adaptive experiments, and network interference require their own
assumptions in place of this baseline. Bounded outcomes, densities, and
smoothness are additional assumptions only where stated.

\textbf{Function classes.}
On $[0,1]^d$, $\mathcal H_d^s(L)$ is the isotropic Hölder ball of radius
$L>0$: put $k=\lceil s\rceil-1$ and $\gamma=s-k\in(0,1]$, bound derivatives
through order $k$, and require the order-$k$ derivatives to be
$\gamma$-Hölder with constant at most $L$. Thus integer orders use a
Lipschitz final derivative convention. The class
$\operatorname{BV}_d(L)$ consists of functions $f\in L^1((0,1)^d)$ whose
distributional derivative is a finite vector measure and for which
$\|f\|_{L^1}+|Df|((0,1)^d)\leq L$.
An $L^\infty$ bound or a particular pointwise version is imposed separately
when needed. Sparse and neural-network classes require a declared
dictionary or architecture and coefficient bounds.

An anisotropic H\"older class with declared block smoothness indices means
coordinatewise H\"older regularity in each block, uniformly in the other
blocks, with all corresponding derivative and H\"older bounds at most
the stated radius. Derivative orders and the integer-order convention in
each block are those just given. Mixed-derivative bounds are additional
restrictions only when explicitly stated.

\textbf{Decision and computational questions.}
Policy regret compares admissible feasible policies under the same law and
constraints. In computational entries, finite data and thresholds have
rational encodings; running time is measured in their total bit length.
Real-valued equilibrium search uses the explicitly declared algebraic or
FIXP encoding. An approximation ratio for maximization is written
$\operatorname{OPT}/\operatorname{ALG}$ and for minimization as
$\operatorname{ALG}/\operatorname{OPT}$, with zero optima and infeasible
instances handled locally. Historical approximation bounds are not
automatically current bounds.
% END ECONOMICS STATEMENT
\end{document}
